\documentclass[11pt,a4paper]{article}
\usepackage[margin=25mm]{geometry}
\usepackage{amsmath,amssymb,booktabs,longtable,array,microtype}
\usepackage[T1]{fontenc}
\usepackage[colorlinks=true,linkcolor=blue,urlcolor=blue]{hyperref}
\setlength{\parindent}{0pt}
\setlength{\parskip}{6pt}
\allowdisplaybreaks
\newcommand{\pos}[1]{\left[#1\right]_{+}}
\newcommand{\off}{\operatorname{off}}
\newcommand{\rn}{\mathcal N}
\DeclareRobustCommand{\key}[1]{\texttt{\detokenize{#1}}}
\pdfstringdefDisableCommands{\let\key\detokenize}
\newcommand{\ind}{\mathbf 1}
\title{BundleFuse: Mathematical Definitions of the Raw Scores}
\author{}
\date{Implementation audit: 6 November 2026}
\begin{document}
\maketitle

Take a candidate item $c$ and a nonempty observed item set $S$, with $c\notin S$.
Each raw score $r_h(c\mid S)$ supplies one number for that candidate.
This document defines these numbers \emph{before} catalogue standardization,
validation-selected weighting, and final interpolation. It covers the eight
dataset configurations and distinguishes their formulas whenever the same score
name has different meanings. All inputs and intermediate quantities used in the
equations are defined here.

\tableofcontents

\section{Shared notation and inputs}

\subsection{Items, bundles, and incidence}
Let $\mathcal I=\{1,\ldots,n\}$ be the item catalogue and
$\mathcal B=\{b_1,\ldots,b_m\}$ the training bundles. Repeated copies of an item
inside one bundle are removed in the graph construction. Define
\[
 B\in\{0,1\}^{m\times n},\qquad B_{bi}=\ind\{i\in b\},\qquad
 x_S\in\{0,1\}^{n},\quad (x_S)_i=\ind\{i\in S\}.
\]
Here $b$ indexes a training bundle, $i$ an arbitrary item, $s,t$ observed items,
and $c$ the candidate. Every formula below uses the unit-weight bundles used by
the dataset pipelines. A separately supplied bundle-weight option in the graph
class is outside this document's scope.

Define the training frequency and training-bundle size by
\[
 f_i=\sum_b B_{bi},\qquad d_b=\sum_iB_{bi}.
\]
Let $\alpha$ and $\gamma$ be the frequency and size exponents, respectively,
and $\beta$ the additional-overlap coefficient. Their default values are
$\alpha=\gamma=0.5$ and $\beta=1.5$. Define
\[
 a_i=(f_i+5)^{-\alpha},\quad v_b=(d_b+1)^{-\gamma},\quad
 D=\operatorname{diag}(a_1,\ldots,a_n),\quad
 V=\operatorname{diag}(v_1,\ldots,v_m).
\]
Thus $D$ is an $n\times n$ diagonal matrix of item-frequency weights; $V$ is an
$m\times m$ diagonal matrix of bundle-size weights. The offsets 5 and 1 are
fixed implementation constants, not counts estimated from the query.

For a scalar $z$, $[z]_+=\max(z,0)$. The operation $\off(A)$ replaces the
diagonal of a square matrix $A$ by zero. For any matrix $F$ of item row vectors,
define row normalization by
\[
 [\rn(F)]_i=\frac{F_i}{\|F_i\|_2+\epsilon_N},\qquad\epsilon_N=10^{-8}.
\]
A zero row remains zero. Norm stabilizers make the normalized length slightly
less than one. A sum over $S$ and a mean over $S$ are kept distinct throughout.

\subsection{Category labels and description arrays}
Let $\ell_i\in\{1,\ldots,q\}$ be the primary label of item $i$.
For POG this is the released clothing category; for Spotify it is the artist;
for dunnhumby it is the product's sub-commodity label. For electronic, food and
clothing it identifies the full parsed category path. Those three datasets also
use a second label $\ell_i^{(2)}$ for the first at most two entries of that path.
Missing retail/commerce categories receive separate item-specific labels.
If explicit primary labels are absent in the general model, they are instead
obtained by clustering the available description arrays (or collaborative
vectors if no description arrays exist).

Let $X^{(j)}\in\mathbb R^{n\times p_j}$ denote the supplied item array for
modality $j$, with $p_j$ coordinates per item. POG provides title, image and
description arrays; Spotify provides an audio array. These are loaded from
\key{title_emb.npy}, \key{content_feature.pt} and
\key{description_feature.pt} as applicable. The scorer consumes these released
numbers; it does not infer their originating encoder from the array values.
Each modality is row-normalized before the constructions below. Retail and
commerce text features are defined explicitly in Section~\ref{sec:text}.

\section{Scores from exact training-item membership}

These scores use item identities and the bundles in which they occur.
First count candidate--seed occurrences; then distinguish normalized direct
support, support from multiple seeds in one historical bundle, and propagated
support.

\subsection{Item co-occurrence: \key{cooc}}
Define
\[
 C=\off(B^\top B),\qquad
 C_{ci}=\begin{cases}\sum_b B_{bc}B_{bi},&c\ne i,\\0,&c=i.\end{cases}
\]
The raw score is
\begin{equation}
 r_{\rm cooc}(c\mid S)=(Cx_S)_c=\sum_{s\in S} C_{cs}.
\end{equation}
Each historical occurrence of a candidate together with a seed contributes
one. A bundle containing several seeds contributes once for each such seed.

\subsection{Frequency-weighted direct score: \key{pair}}
For each training bundle define its weighted overlap and overlap count:
\[
 o_b(S)=\sum_{s\in S}B_{bs}a_s,\qquad
 k_b(S)=\sum_{s\in S}B_{bs}=|b\cap S|.
\]
The direct score is
\begin{equation}\label{eq:pair}
 p_c(S)=r_{\rm pair}(c\mid S)
 =\sum_b B_{bc}v_bo_b(S)
 =\bigl[B^\top VBDx_S\bigr]_c.
\end{equation}
Relative to ordinary co-occurrence, rarer seeds have larger $a_s$, and larger
historical bundles have smaller $v_b$. There is no candidate-frequency factor
$a_c$ in this direct score.

The stored zero-diagonal pair matrix is
\[
 W=\off(B^\top VBD),\qquad
 W_{cs}=a_s\sum_b B_{bc}v_bB_{bs}\quad(c\ne s).
\]
For eligible candidates $c\notin S$, $(Wx_S)_c=p_c(S)$. Nevertheless, the
implementation computes the full vector $p=B^\top VBDx_S$, \emph{including}
its seed coordinates, before propagation. Substituting $Wx_S$ for $p$ inside
the propagated score would therefore change that score. Generally $W$ is not
symmetric because only its seed-side index receives $a_s$.

\subsection{Additional joint-overlap score: \key{delta_inc}}
The extra contribution is
\begin{equation}\label{eq:delta}
 \delta_c(S)=\beta\sum_b B_{bc}v_bo_b(S)[k_b(S)-1]_+.
\end{equation}
Its matrix form is
\[
 \delta(S)=\beta B^\top V
 \left((BDx_S)\odot[Bx_S-\ind]_+\right),
\]
where $\odot$ denotes coordinatewise multiplication and $\ind$ is the
$m$-dimensional all-ones vector. A training bundle with only one seed adds
no extra contribution. A bundle with two seeds multiplies its direct
contribution by an additional $\beta$; a bundle with three adds $2\beta$ times
that contribution. This is additional credit for \emph{joint membership in
the same historical bundle}, not merely for two independent pair counts.

More explicitly, define the candidate--seed--seed count
\[
 H_{cst}=\sum_b v_bB_{bc}B_{bs}B_{bt}.
\]
Then
\begin{equation}
 \delta_c(S)=\beta\sum_{\{s,t\}\subseteq S}(a_s+a_t)H_{cst},
\end{equation}
where each unordered pair of distinct seeds is included once. This identity
also shows that $\delta_c(S)=0$ for a single-seed query.

\subsection{Direct plus additional joint overlap: \key{simp}}
The complete incidence score is
\begin{equation}
 g_c(S)=r_{\rm simp}(c\mid S)=p_c(S)+\delta_c(S).
\end{equation}
The symbols $g_c$ and $\delta_c$ denote different quantities: the former
includes direct support; the latter contains only the extra contribution.
The commerce profile replaces $p$ by an augmented direct score defined in
Section~\ref{sec:commerce}.

\subsection{Propagated scores: \key{diff2_pair} and \key{diff2}}
Define the $n\times n$ operator
\begin{equation}\label{eq:L}
 L=\off\!\left(DB^\top EBD\right),\qquad
 E=\operatorname{diag}\bigl((d_b+1)^{-1}\bigr)_{b=1}^m.
\end{equation}
In particular,
\[
 L_{ci}=a_ca_i\sum_b\frac{B_{bc}B_{bi}}{d_b+1}\ (c\ne i),\qquad L_{cc}=0.
\]
Here both item endpoints receive a frequency weight. $L$ is symmetric and
nonnegative, but its rows need not sum to one. It is a normalized adjacency
operator, not a graph Laplacian.

The raw propagated scores are
\begin{align}
 r_{\rm diff2\_pair}(c\mid S)&=(Lp(S))_c,
 \\
 r_{\rm diff2}(c\mid S)&=(Lg(S))_c
 =(Lp(S))_c+(L\delta(S))_c.
\end{align}
An intermediate item $i$ passes its score $p_i$ or $g_i$ to candidate $c$
with weight $L_{ci}$. All items, including seeds, participate as intermediates.
Because the normalizations in $p$ and $L$ differ, the first expression is not
generally $L^2x_S$. The commerce profile uses a different propagation matrix.

\subsection{Popularity: \key{pop}}
\begin{equation}
 r_{\rm pop}(c\mid S)=\log(1+f_c).
\end{equation}
It depends on candidate training frequency only and is identical for all queries.

\section{Association strengths and their compressed profiles}

Co-occurrence counts favor frequent items. The next score compares a pair
count with a product of marginal counts. SVD then compresses an item's entire
association row, so candidates can be compared through similar patterns of
association with other items.

\subsection{Positive item association: \key{ppmi}}
Using the zero-diagonal $C$, define
\[
 q_i=\sum_jC_{ij},\qquad Q=\sum_{i,j}C_{ij}.
\]
The endpoint count $q_i$ differs from bundle frequency $f_i$:
$q_i=\sum_bB_{bi}(d_b-1)$. The symmetric matrix $M$ is
\begin{equation}
 M_{ij}=\begin{cases}
 \displaystyle\pos{\log\!\left(
 \frac{C_{ij}Q}{(q_i+10^{-9})(q_j+10^{-9})}+10^{-12}\right)},&C_{ij}>0,\\
 0,&C_{ij}=0.
 \end{cases}
\end{equation}
The name PPMI denotes positive pointwise mutual information. The raw score is
\begin{equation}
 r_{\rm ppmi}(c\mid S)=(Mx_S)_c=\sum_{s\in S}M_{cs}.
\end{equation}
Only pair counts larger than their marginal-based reference receive positive
credit. The logarithm is natural.

\subsection{Compressed association-profile similarity: \key{csvd}}
Let $r=\min(r_0,\max(1,n-2))$, where $r_0=128$ by default. A truncated
singular value decomposition supplies right singular directions
$V_r\in\mathbb R^{n\times r}$, left singular directions
$U_r\in\mathbb R^{n\times r}$ and the diagonal matrix of retained singular
values $\Sigma_r\in\mathbb R^{r\times r}$. Define
\begin{equation}
 Z=MV_r\simeq U_r\Sigma_r,\qquad U=\rn(Z).
\end{equation}
The implementation uses a randomized truncated decomposition with random
state zero; $MV_r$ is the operational definition of the returned feature
matrix. If no associations are stored, it uses zero features.
The raw score is
\begin{equation}
 r_{\rm csvd}(c\mid S)=U_c\left(\frac1{|S|}\sum_{s\in S}U_s\right)^\top
 =\frac1{|S|}\sum_{s\in S}U_cU_s^\top.
\end{equation}
This averages cosine similarities of compressed association profiles, up to
the norm stabilizer. It is not the entry $(c,s)$ of the reconstructed matrix
$U_r\Sigma_rV_r^\top$. Before row normalization the profile Gram matrix is
$U_r\Sigma_r^2U_r^\top$.

Even though $M$ is symmetric, it need not be positive semidefinite. SVD retains
directions by singular-value magnitude, not by positive eigenvalues alone.
Orthogonality concerns the \emph{columns} of singular-vector matrices;
individual item \emph{rows} of $U_r\Sigma_r$ do not have equal norms.
This is why row normalization is a separate operation.

\section{Scores from primary labels and soft feature groups}

These constructions share information between different item identities.
The hard-label score uses the supplied category or artist; the soft-label
score uses clusters of item features. These are different group definitions.

\subsection{Counts of labels inside training bundles}
Define the bundle--label count matrix
\[
 T\in\mathbb R^{m\times q},\qquad
 T_{ba}=\sum_iB_{bi}\ind\{\ell_i=a\}.
\]
Unlike $B$, $T$ may contain values greater than one. Let
\[
 h_a=\sum_bT_{ba},\qquad H=\sum_a(h_a+10^{-9}),\qquad A=T^\top T.
\]
Here $A\in\mathbb R^{q\times q}$ counts label pairs through
\[
 A_{at}=\sum_{b=1}^m T_{ba}T_{bt}.
\]
For example, a training bundle containing two items of label $a$ and one item
of a different label $t$ contributes $2\times1=2$ to $A_{at}$.
Thus this is a count of item combinations across labels, rather than just a
count of bundles in which both labels are present.

The symbol $J\in\mathbb R_{\ge0}^{q\times q}$ denotes the label-pair count
matrix supplied to the association transform below. It is constructed from
$T$, with the diagonal choice determined by the score:
\[
 J=\begin{cases}
 \off(T^\top T),&\text{different-label association used by \key{cat}},\\
 T^\top T,&\text{diagonal-retaining commerce \key{homo} association}.
 \end{cases}
\]
For distinct labels $a\ne t$, both choices have
$J_{at}=\sum_bT_{ba}T_{bt}$. For $a=t$, the first choice sets $J_{aa}=0$;
the second retains $J_{aa}=\sum_bT_{ba}^2$. The retained square includes
self-item contributions; it is not a count restricted to distinct item pairs.
Therefore $J$ is not a new dataset, a learned matrix, or a query-dependent
quantity. It is the training label-count product, with or without its diagonal.

For either choice of $J$, define the transform using the marginals above:
\begin{equation}\label{eq:catpmi}
 [\mathcal P_T(J)]_{at}=
 \begin{cases}
 \displaystyle\pos{\log\!\left(\frac{J_{at}H}
 {(h_a+10^{-9})(h_t+10^{-9})}+10^{-12}\right)},&J_{at}>0,\\
 0,&J_{at}=0.
 \end{cases}
\end{equation}
The label marginals are occurrence counts in $T$, not row sums of $J$.
Consequently this transform and the item PPMI construction use different
marginals. Their shared logarithmic form does not make their definitions
interchangeable.

\subsection{Different-label association and slot adjustment: \key{cat}}
Set $P^{\rm cat}=\mathcal P_T(\off(A))$. Its diagonal is zero.
Define the unadjusted category score
\[
 r_{\rm cat}^{\rm base}(c\mid S)=\sum_{s\in S}P^{\rm cat}_{\ell_s,\ell_c}.
\]
For the general non-music path, define for each label $a$
\[
 n_a^{(1)}=\sum_b\ind\{T_{ba}\ge1\},\qquad
 n_a^{(2)}=\sum_b\ind\{T_{ba}\ge2\},\qquad
 \eta_a=\frac{n_a^{(2)}+0.5}{n_a^{(1)}+1}.
\]
The slot adjustment is
\[
 r_{\rm slot}(c\mid S)=-\sum_{s\in S}
 \ind\{\ell_c=\ell_s\}\ind\{\eta_{\ell_s}<0.15\}(1-\eta_{\ell_s}).
\]
Then
\begin{equation}
 r_{\rm cat}(c\mid S)=
 \begin{cases}
 r_{\rm cat}^{\rm base}(c\mid S)+0.5r_{\rm slot}(c\mid S),&\text{general non-music},\\
 r_{\rm cat}^{\rm base}(c\mid S),&\text{music or research retail/commerce}.
 \end{cases}
\end{equation}
Thus a category seldom repeated within training bundles receives a negative
same-category adjustment in the first case. Multiple seeds with that category
each contribute the adjustment. Retail/commerce research scores do not apply it.

\subsection{Shared-label score: \key{homo}}
\paragraph{Spotify and Spotify sparse.}
Let $a(i)$ and $b(i)$ be item $i$'s artist and album identifiers, and let
$N_a$ and $N_b$ denote the numbers of catalogue items with those identifiers.
The raw score is
\begin{equation}
 r_{\rm homo}(c\mid S)=\sum_{s\in S}
 \left(\frac{\ind\{a(c)=a(s)\}}{N_{a(c)}+5}
 +2\frac{\ind\{b(c)=b(s)\}}{N_{b(c)}+2}\right).
\end{equation}
If album labels are unavailable, omit the album term. These denominators count
catalogue items, not training occurrences.

\paragraph{General model with supplied hard group labels.}
For group type $j$, let $g_j(i)$ be the supplied label, $N_{ja}$ the number of
items in group $a$, and $\omega_j$ its supplied weight. Then
\[
 r_{\rm homo}(c\mid S)=\sum_j\omega_j\sum_{s\in S}
 \frac{\ind\{g_j(c)=g_j(s)\}}{N_{j,g_j(c)}+1}.
\]
These groups are supplied labels; they need not be clusters of features.

\paragraph{General non-music model without supplied groups.}
The fallback definition retains the diagonal of the label count product:
\[
 r_{\rm homo}(c\mid S)=\sum_{s\in S}(T^\top T)_{\ell_s,\ell_c}.
\]
This is a raw count score, with no logarithmic transform. It is available in
the extended profile but outside the POG research selection menu.
The commerce research version is defined in Section~\ref{sec:commerce}.

\subsection{Soft feature-group association: \key{softcat}}
Let $F\in\mathbb R^{n\times p}$ be the features used for clustering.
For POG's research profile, $F$ is the row-normalized concatenation of the
normalized title, image and description arrays that are present. Apply
MiniBatchKMeans to the item rows, with
$k=\max(2,\min(k_0,n-1))$, $k_0=400$ by default, random state zero,
three initializations and batch size 4096. Let $\mu_a$ be cluster $a$'s
arithmetic centre and $\widehat\mu_a=\mu_a/(\|\mu_a\|_2+10^{-8})$.
A feature group here is precisely one of these clusters; it is not a
training bundle and need not coincide with a clothing category.

With temperature $\tau_s=0.05$, define soft membership
\begin{equation}
 A^{\rm soft}_{ia}=
 \frac{\exp(F_i\widehat\mu_a^\top/\tau_s-m_i)}
 {\sum_{t=1}^k\exp(F_i\widehat\mu_t^\top/\tau_s-m_i)+10^{-9}},\qquad
 m_i=\max_t F_i\widehat\mu_t^\top/\tau_s.
\end{equation}
Each item is assigned a weight for every group rather than only one hard label.
Define a separate matrix, $J^{\rm soft}\in\mathbb R^{m\times k}$, by
\[
 J^{\rm soft}=BA^{\rm soft},\qquad
 J^{\rm soft}_{ba}=\sum_{i=1}^nB_{bi}A^{\rm soft}_{ia}
                  =\sum_{i\in b}A^{\rm soft}_{ia}.
\]
Row $b$ describes one training bundle and column $a$ one feature cluster.
Its entry is the sum of the membership weights in cluster $a$ of every item
in bundle $b$. For example, membership weights $0.8$, $0.3$ and $0.1$ for
three items in that bundle give $J^{\rm soft}_{ba}=1.2$.
It is a soft group count, not a probability, and can exceed one.
This $m\times k$ matrix is distinct from the $q\times q$ hard-label pair
matrix $J$ defined in Section~4.1.

The product $((J^{\rm soft})^\top J^{\rm soft})_{at}
=\sum_bJ^{\rm soft}_{ba}J^{\rm soft}_{bt}$ measures how much groups $a$ and
$t$ are jointly represented across training bundles. Expanding each product
also counts a single item's membership in both groups. Remove those
same-item contributions, then remove the group diagonal:
\begin{equation}
 C^{\rm soft}=\pos{\off\!\left(
 (J^{\rm soft})^\top J^{\rm soft}
 -\sum_i f_i(A^{\rm soft}_i)^\top A^{\rm soft}_i\right)}.
\end{equation}
For distinct groups $a,t$, this equals
$\sum_b\sum_{i\ne j}B_{bi}B_{bj}A^{\rm soft}_{ia}A^{\rm soft}_{jt}$,
apart from numerical clipping. The diagonal is then excluded by construction.
Define \emph{catalogue} group masses $u_a=\sum_iA^{\rm soft}_{ia}$ and
$U_0=\sum_au_a$, and set
\begin{equation}
 P^{\rm soft}_{at}=\pos{\log\!\left(
 \frac{C^{\rm soft}_{at}U_0}{u_au_t+10^{-9}}+10^{-12}\right)}.
\end{equation}
The score is
\begin{equation}
 r_{\rm softcat}(c\mid S)=A^{\rm soft}_cP^{\rm soft}
 \left(\sum_{s\in S}A^{\rm soft}_s\right)^\top.
\end{equation}
The group co-occurrences are training-bundle counts while their denominator
uses catalogue group masses. This is the implemented association transform;
it should not be silently replaced by a probability model with different
marginals. This score is constructed for the non-music general path when
description features are available.

\section{Description, residual-description, and bridge scores}

All dense-feature scores have one scoring rule. If $F$ is the relevant
$n\times p$ feature matrix, define
\begin{equation}\label{eq:feature}
 r_F(c\mid S)=F_c\left(\frac1{|S|}\sum_{s\in S}F_s\right)^\top.
\end{equation}
The observed mean is not normalized again. Different score names specify
different constructions of $F$.

\subsection{Direct description similarity: \key{content}}
For the general model,
\[
 F^{\rm content}=\rn\!\left([X^{(1)}\mid\cdots\mid X^{(J)}]\right),\qquad
 r_{\rm content}=r_{F^{\rm content}}.
\]
The vertical separators mean concatenate coordinates, not sum modalities.
For research retail/commerce, the content matrix is the normalized
128-dimensional text representation in Section~\ref{sec:text}.

\subsection{Within-group residual descriptions}
Define $g(i)$ to be the label used for centring and
$\mathcal G_a=\{i:g(i)=a\}$. For modality $j$, compute
\begin{equation}
 \overline X_a^{(j)}=\frac1{|\mathcal G_a|}\sum_{i\in\mathcal G_a}X_i^{(j)},
 \qquad R_i^{(j)}=X_i^{(j)}-\overline X_{g(i)}^{(j)},
 \qquad F^{(j)}=\rn(R^{(j)}).
\end{equation}
The centre is this coordinatewise arithmetic mean over exactly the specified
item group, not a mean over all features or over the query.
For non-music items, $g(i)=\ell_i$: the supplied clothing category in POG.
For music, $g(i)$ comes from a separate MiniBatchKMeans fit on the provided
clustering features (concatenated audio features in the released data).
Music centring therefore uses feature clusters, not artists.

\begin{center}
\begin{tabular}{lll}\toprule
Raw score key & Modality $j$ & Raw value\\\midrule
\key{style} & title/general content & $r_{F^{(j)}}(c\mid S)$\\
\key{imgstyle} & image & $r_{F^{(j)}}(c\mid S)$\\
\key{descstyle} & description & $r_{F^{(j)}}(c\mid S)$\\
\key{audiostyle} & audio & $r_{F^{(j)}}(c\mid S)$\\\bottomrule
\end{tabular}
\end{center}
Subtracting the group mean removes the shared mean description before comparing
the remaining coordinates. It does not prove statistical independence from
the group, and row normalization need not preserve a zero group mean.
The POG research menu selects \key{style} and \key{imgstyle}; description
residuals can still enter its bridge even though \key{descstyle} is not a
standalone menu entry.

\subsection{History-propagated description similarity: \key{bridge}}
In the general non-music path define
\[
 F^0=\rn([F^{(1)}\mid\cdots\mid F^{(J)}]).
\]
For the Spotify research profile, $F^0=F^{\rm audio}$.
Using $L$ from Equation~\eqref{eq:L}, construct
\begin{equation}
 N=\rn(LF^0),\qquad
 F^{\rm bridge}=\rn(F^0+\lambda_BN),\qquad
 r_{\rm bridge}=r_{F^{\rm bridge}},\quad\lambda_B=0.5.
\end{equation}
The row $(LF^0)_c=\sum_iL_{ci}F^0_i$ is the sum of descriptions of other
items, weighted by their normalized historical co-membership with candidate
$c$. Mixing it with $F^0_c$ gives the bridge vector. Both the candidate and
each seed use these mixed vectors.

For research dunnhumby, start from normalized text content $F^{\rm content}$
instead of category-centred residuals:
\[
 F^{\rm bridge}=\rn\!\left(F^{\rm content}
 +0.5\rn(LF^{\rm content})\right).
\]
The commerce bridge has a separate sparse-text formula below.

\subsection{Text features and the \key{tfidf} score}\label{sec:text}
For dunnhumby concatenate department, commodity, sub-commodity and brand text.
For commerce concatenate item title and its parsed category path. Tokenization
uses the vectorizer's default word tokens, lowercasing, English stop-word
removal, unigrams and bigrams. Retain at most 8000 terms for dunnhumby and
12000 for the commerce research profile. Vocabulary truncation is based on
corpus term frequency after the vectorizer's filtering.

Let $t_{iw}$ be term $w$'s count in item $i$ and
$\mathrm{df}_w=\sum_i\ind\{t_{iw}>0\}$. The pre-normalization entry is
\[
 \widetilde H_{iw}=
 \begin{cases}
 (1+\log t_{iw})\left(1+\log\frac{1+n}{1+\mathrm{df}_w}\right),&t_{iw}>0,\\
 0,&t_{iw}=0.
 \end{cases}
\]
Let $H$ be its row-unit-normalized sparse matrix, with zero rows kept zero.
The TF-IDF score is
\begin{equation}
 r_{\rm tfidf}(c\mid S)=(HH^\top x_S)_c=\sum_{s\in S}H_cH_s^\top.
\end{equation}
This uses a \emph{sum}, not the mean in Equation~\eqref{eq:feature}.
For the dense content matrix used by the research basket profiles, compute
\[
 F^{\rm content}=\rn(HV_{r_T}),\qquad r_T=\min(128,\#\text{terms}-1),
\]
where $V_{r_T}$ are the retained right singular directions of $H$.
Vocabulary and document-frequency counts use the catalogue item texts;
they do not use hidden query labels. The retail research menu does not select
\key{tfidf}, although its scorer makes this score available as an auxiliary.

\section{Candidate support across multiple seeds: \key{rao}}

This score starts from candidate--seed affinities and then multiplies support
from distinct seeds. It is different from the historical triple count in
$\delta$: it can be nonzero when candidate--seed pairs occurred in different
training bundles.

Let $F^R$ be the row-normalized features used in this score. They are title
residuals for POG, audio residuals for Spotify, dense bridge features for
research dunnhumby, and dense commerce bridge features for research commerce.
Let $\zeta=0.15$, except $\zeta=0.25$ for research commerce. Define
\begin{equation}
 w_c(s)=W_{cs}+\zeta\pos{F^R_c(F^R_s)^\top},\qquad
 A_c(S)=\sum_{s\in S}w_c(s).
\end{equation}
Here $W$ is the zero-diagonal direct-pair matrix already defined. The additional
term is positive description similarity; it is not a validation-fitted
candidate--seed weight.

When both descriptions and primary labels are available, define
\[
 G_{st}=0.5\ind\{\ell_s=\ell_t\}
       +0.5\pos{F^R_s(F^R_t)^\top}\quad(s\ne t),\qquad G_{ss}=1.
\]
With only descriptions, use their positive dot product; with only labels, use
the equal-label indicator. With neither, set $G$ to the identity matrix.
The seed-pair coefficient is
\begin{equation}
 K_{st}=\begin{cases}
 0,&s=t,\\
 [1+\gamma_R G_{st}]_+,&s\ne t,\ \gamma_R<0,\\
 1+\gamma_R G_{st},&s\ne t,\ \gamma_R\ge0.
 \end{cases}
\end{equation}
The dataset rule is $\gamma_R=1$ for Spotify variants and $\gamma_R=-1$ for
POG, retail and commerce. The raw score is
\begin{equation}
 r_{\rm rao}(c\mid S)=
 \frac{\displaystyle\sum_{s\in S}\sum_{t\in S,\ t\ne s}
 w_c(s)K_{st}w_c(t)}{A_c(S)+0.05}.
\end{equation}
The numerator counts ordered pairs, so a symmetric pair contributes twice.
For POG, two similar seeds have smaller $K_{st}$ and supply less extra
credit; this does not subtract from their ordinary direct scores.
For music the sign is reversed: similar seeds have larger $K_{st}$.
Accordingly the music score should not be described as a diversity penalty.
A single seed always gives zero.

The graph class also exposes a simpler auxiliary operator using $w_c(s)=W_{cs}$
and $K_{st}=1$ for $s\ne t$:
\begin{equation}
 r_{e_2}(c\mid S)=\frac{
 \left[\left(\sum_sw_c(s)\right)^2-\sum_sw_c(s)^2\right]_+}
 {\sum_sw_c(s)+0.05}.
\end{equation}
It is not a separate default selection-menu score. The legacy combined graph
helper can add $0.35r_{e_2}$ to $p+\delta$; the main BundleFuse scoring path
does not make that addition implicitly.

\section{Household, user, session, and commerce variants}

\subsection{Retail household co-occurrence: \key{hh_cooc}}
Let $Y_{ui}$ be the number of training baskets of household $u$ that contain
item $i$. Counts are retained across baskets. Define
\[
 d_u^Y=\sum_iY_{ui},\qquad
 C^{\rm hh}=\off\!\left(Y^\top\operatorname{diag}((d_u^Y+1)^{-1})Y\right).
\]
The research dunnhumby score is
\begin{equation}
 r_{\rm hh\_cooc}(c\mid S)=(C^{\rm hh}Dx_S)_c.
\end{equation}
The household normalization reduces contributions from households with larger
total histories. The seed side receives the same $a_s$ weight as direct
bundle support. Same-household items can contribute even when they occurred
in different baskets.

\subsection{Commerce research profile}\label{sec:commerce}
This subsection applies to electronic, food and clothing with their default
research profile. It modifies several shared scores.

\paragraph{User, session, and pretraining matrices.}
Let $Y^u_{ui}$ count occurrences of item $i$ across the training bundles
assigned to user $u$, and $Y^s_{vi}$ count occurrences across training bundles
assigned to session $v$. Ignore bundles without a mapped user/session for
that corresponding matrix. Define
\[
 C^u=\off((Y^u)^\top Y^u),\qquad C^s=\off((Y^s)^\top Y^s).
\]
Let $Y^{\rm pre}_{ui}$ count rows of the supplied user--item pretraining file
for user $u$, item $i$. Its third column is ignored. Let
$d_u^{\rm pre}=\sum_iY^{\rm pre}_{ui}$. Define
\[
 C^{\rm pre}=\off\!\left((Y^{\rm pre})^\top
 \operatorname{diag}\left(\max(d_u^{\rm pre},1)^{-1}\right)Y^{\rm pre}\right).
\]
Then the raw \key{pretrain} score is
\begin{equation}
 r_{\rm pretrain}(c\mid S)=(C^{\rm pre}x_S)_c.
\end{equation}

\paragraph{Augmented incidence and propagation.}
Using the bundle-only direct score $p$ in Equation~\eqref{eq:pair}, set
\begin{align}
 p^{\rm com}&=p+0.35C^ux_S+0.35C^sx_S,\\
 g^{\rm com}&=p^{\rm com}+\delta,\\
 D_1&=\operatorname{diag}((f_i+1)^{-1/2}),\\
 P^{\rm com}&=D_1(C+0.5C^s+0.5C^u+0.2C^{\rm pre})D_1.
\end{align}
Thus $r_{\rm pair}=p^{\rm com}$, $r_{\rm simp}=g^{\rm com}$,
$r_{\rm diff2\_pair}=P^{\rm com}p^{\rm com}$ and
$r_{\rm diff2}=P^{\rm com}g^{\rm com}$, evaluated at candidate coordinate $c$.
The extra $\delta$ still comes only from training bundles. $P^{\rm com}$ is
not row-normalized. Item \key{cooc}, \key{ppmi} and \key{csvd} continue to use
the bundle-only graph, not $C^u+C^s+C^{\rm pre}$.

\paragraph{Two-level category score.}
Construct $T^{(1)}$ from full-path labels and $T^{(2)}$ from the first at most
two path entries. Apply Equation~\eqref{eq:catpmi} separately with each
matrix's own marginals. Define
\[
 P^{(j)}=\mathcal P_{T^{(j)}}((T^{(j)})^\top T^{(j)}),\quad j\in\{1,2\}.
\]
Here the diagonal is \emph{retained}. The commerce \key{homo} score is
\begin{equation}
 r_{\rm homo}(c\mid S)=\sum_{s\in S}
 \left(P^{(1)}_{\ell_s,\ell_c}
 +0.5P^{(2)}_{\ell_s^{(2)},\ell_c^{(2)}}\right).
\end{equation}
The separate \key{cat} score uses the off-diagonal transform of $T^{(1)}$
and no slot adjustment.

\paragraph{Sparse text bridge and dense features for concordance.}
With sparse TF-IDF $H$ and dense normalized text features $F^{\rm content}$,
define
\begin{align}
 H^{\rm bridge}&=\rn_0\!\left(H+1.5P^{\rm com}H
                          +0.5(P^{\rm com})^2H\right),\\
 F^{R}&=\rn\!\left(F^{\rm content}+1.5P^{\rm com}F^{\rm content}
                          +0.5(P^{\rm com})^2F^{\rm content}\right).
\end{align}
Here $\rn_0$ is exact nonzero-row unit normalization with zero rows preserved,
as used by the sparse normalization routine. The commerce raw bridge score is
\begin{equation}
 r_{\rm bridge}(c\mid S)=
 \bigl[H^{\rm bridge}(H^{\rm bridge})^\top x_S\bigr]_c.
\end{equation}
This is a sum of sparse-text similarities, not a mean of dense-vector
similarities. $F^R$ is used for \key{rao}, not for this bridge score.

\subsection{General auxiliary incidence scores}
For an explicitly supplied context-by-item incidence matrix
$I\in\mathbb R^{t\times n}$, the general model supports a named score
\begin{equation}
 r_{\rm aux}(c\mid S)=
 \begin{cases}
 [I^\top IDx_S]_c,&\text{seed-frequency weighting enabled},\\
 [I^\top Ix_S]_c,&\text{otherwise}.
 \end{cases}
\end{equation}
The matrix may already contain context-degree normalization. The general
path does not remove the diagonal of $I^\top I$; eligible candidates exclude
the seeds. The loader's extended retail \key{hh_cooc} and commerce
\key{pretrain} incidences use $I_{ui}=Y_{ui}/\sqrt{\sum_jY_{uj}+1}$.
They are not identical constructions to the research commerce pretraining
normalization above.

\subsection{Personalized auxiliaries: \key{user_freq}, \key{user_cooc}}
These are available in the general extended path when a user history is
supplied. For query user $u$, let $n_{ui}$ be the number of occurrences of
item $i$ in that supplied history, and $h_{ui}=\ind\{n_{ui}>0\}$. Then
\begin{align}
 r_{\rm user\_freq}(c\mid S,u)&=n_{uc},\\
 r_{\rm user\_cooc}(c\mid S,u)&=(Ch_u)_c.
\end{align}
The first retains repetitions; the second uses the distinct-history indicator.
Neither uses the current query seeds directly. Unknown or absent users receive
zero scores. These are not selected by the default research basket menus.

\section{Inventory and the boundary of this document}

\subsection{Default research selection menus}
The following lists identify which raw sources are offered for ordinary
validation weight selection. Offered does not mean selected with positive
weight in a fitted model.

\begin{longtable}{>{\raggedright\arraybackslash}p{0.22\textwidth}>{\raggedright\arraybackslash}p{0.70\textwidth}}
\toprule Dataset & Raw source menu\\\midrule\endhead
POG, POG dense &
\key{cat}, \key{softcat}, \key{style}, \key{imgstyle}, \key{bridge},
\key{simp}, \key{diff2}, \key{pop}, \key{cooc}, \key{ppmi}, \key{csvd}.\\[4pt]
Spotify, Spotify sparse &
\key{simp}, \key{diff2}, \key{pop}, \key{cooc}, \key{ppmi}, \key{csvd},
\key{cat}, \key{homo}, \key{audiostyle}, \key{bridge}.\\[4pt]
Dunnhumby &
\key{simp}, \key{diff2}, \key{pop}, \key{cooc}, \key{hh_cooc},
\key{ppmi}, \key{csvd}, \key{cat}, \key{content}, \key{bridge}.\\[4pt]
Electronic, food, clothing &
\key{simp}, \key{diff2}, \key{cooc}, \key{ppmi}, \key{csvd},
\key{cat}, \key{homo}, \key{tfidf}, \key{bridge}, \key{pretrain}.\\
\bottomrule
\end{longtable}
The fitting/scoring logic additionally uses \key{pair}, \key{diff2_pair}
and \key{rao} to form controlled combinations with these sources.
\key{delta_inc} is directly addressable in the selected general scorer;
the ordinary full scorer uses it internally through \key{simp} and
\key{diff2}. Placeholder entries \key{contrast} and \key{img} in the basket
reference scorer contain no numerical score and are not menu entries.

\subsection{Optional content compression}
The default configurations set \key{content_rank} to null. If optional
content compression is enabled with rank $r_C$, the stored content, style
and bridge matrices with more than $r_C$ columns are replaced by
$FV_{r_C}$ from a truncated SVD. They are not row-normalized again.
Equation~\eqref{eq:feature} then uses these projected matrices; calling the
result an exact cosine similarity would be incorrect. The research basket
profiles reject this override to preserve their sparse/dense constructions.

\subsection{Experimental comparison operator, outside BundleFuse}
The novelty experiments contain an amplified-neighbour comparator.
It is not an additional BundleFuse raw source. To distinguish it precisely,
define
\[
 t_b(S)=\frac{\sum_{s\in S}B_{bs}a_s}{\sqrt{d_b}\sqrt{|S|}}.
\]
Keep all bundles or only the $k_N$ largest $t_b$, using bundle index to break
ties; let $\mathcal K(S)$ denote those retained indices. For chosen power
$p_N$, the comparator is
\[
 r_{\rm amp}(c\mid S)=\sum_{b\in\mathcal K(S)}B_{bc}\,t_b(S)^{p_N}.
\]
This is a predecessor-inspired control; its implementation does not claim an
exact reproduction of another published method.

\subsection{What happens after a raw score}
The scalar functions above are the raw inputs to subsequent standardization
and fusion. Their fixed construction constants (such as 5, 0.5 and 0.05)
are distinct from validation-selected fusion weights. Selecting a zero fusion
weight does not change a raw source's definition, but serving can omit its
computation when no remaining selected score depends on it.

\section*{Implementation trace for verification}
This document was derived from executable expressions in the following files,
rather than earlier mathematical documents. The mapping is provided only for
checking the equations; no external document is required to define a symbol.
\begin{itemize}
\item \key{bundlefuse/operators.py}: graph, incidence surplus, propagation,
item PPMI/SVD, category transform, soft groups, residuals and concordance.
\item \key{bundlefuse/model.py} and \key{bundlefuse/selected_signals.py}:
construction of description matrices and assignment of raw score names.
\item \key{bundlefuse/basket_signals.py}: research retail and commerce variants.
\item \key{bundlefuse/datasets/base.py} and \key{bundlefuse/datasets/baskets.py}:
label sources, text arrays and auxiliary incidence construction.
\item \key{experiments/novelty/raw.py}: amplified-neighbour comparison only.
\end{itemize}
The experimental source files and running experiment configuration were not
modified to prepare this document.

For numerical verification, 87 comparisons on a small independent synthetic
bundle collection checked the displayed core equations and the full versus
selected scoring paths for fashion, music, retail and commerce, including
single-seed and multiple-seed queries. The comparisons passed with relative
tolerance $3\times10^{-5}$ and absolute tolerance $3\times10^{-6}$.
\end{document}
